\subsection{Funções comuns de teoria dos números}

\subsubsection{Função totiente de Euler}

Seja $\phi(n)$ a função totiente de euler.

Código para calcular $\phi(n)$ em $O(\sqrt{n})$:

\inccode{algtn/phi.cpp}

Código para calcular $\phi(n)$ de $1$ a $n$ em $O(n \log \log n)$ com crivo:

\inccode{algtn/phi_sieve.cpp}

Fatos conhecidos:
$$\phi(p_1^{\alpha_1} \dots p_k^{\alpha_k}) = p_1^{\alpha_1} \dots p_k^{\alpha_k} \left(1 - \dfrac{1}{p_1} \right) \dots \left(1 - \dfrac{1}{p_k} \right) $$
$$\sum\limits_{d | n} \phi(d) = n$$
Para calcular $x^n$ mod $m$ rápido (mesmo que $\gcd(x, m) \neq 1$) podemos usar:
$$x^n \equiv x^{\phi(m) + [n \mod \phi(m)]} ( \mod m)$$

\subsubsection{Funções quantidade e soma de divisores}

Seja $n = p_1^{e_1} \dots p_k^{e_k}$. Então
$$d(n) = \prod\limits_{i=1}^k (e_i+1)$$
$$\sigma(n) = \prod\limits_{i=1}^k \dfrac{p_i^{e_i+1}-1}{p_i-1}$$

É possível calculá-las em $O(\sqrt{n})$ ingenuamente. Mas se você usar o crivo super-poderoso e ir subindo no vetor de primos, dá até pra calcular a soma dos divisores de vários $n \leq 10^{16}$. (dá pra usar crivo normal também, ele só não vai chegar a $10^{16}$).

\subsubsection{Função de mobius}

$\mu(n) = 0$ se $p^2 | n$ para algum $p$. Caso contrário, $\mu(n) = (-1)^k$ em que $k$ é o número de fatores primos distintos de $n$. Por exemplo, $\mu(1) = 1$, $\mu(2) = -1$, $\mu(3) = -1$, $\mu(4) =0$, $\mu(6) = 1$. Calcula com crivo:

\inccode{algtn/mobius_sieve.cpp}

\subsubsection{Funções aritméticas multiplicativas}

Funções nos inteiros tais que $f(ab) = f(a)f(b)$ quando $\gcd(a, b) = 1$. Na grande maioria dos casos, é fácil adaptar o crivo acima. O único problema é caso fique difícil de fazer o caso em que o primo divide $i$. O que você pode fazer é criar uma array $cnt$ indicando o expoente do menor fator primo. Aqui um exemplo em que $f(p^k) = k$:

\inccode{algtn/mult_sieve.cpp}